Trong những năm gần đây, nhu cầu chăm sóc và theo dõi sức khỏe thú cưng ngày càng được quan tâm. Cùng với sự phát triển của các cơ sở thú y, lượng thông tin cần quản lý như thông tin khách hàng, thú cưng, lịch hẹn, bác sĩ thú y, hồ sơ khám bệnh, hóa đơn và phòng chăm sóc cũng ngày càng tăng. Việc quản lý các thông tin này một cách thủ công có thể gây khó khăn trong việc lưu trữ, tra cứu, cập nhật và đảm bảo tính chính xác của dữ liệu.

Xuất phát từ nhu cầu trên, đề tài  Chăm sóc Sức khỏe Thú cưng được thực hiện nhằm phân tích và xây dựng một cơ sở dữ liệu phục vụ việc quản lý các thông tin và nghiệp vụ cơ bản của một cơ sở chăm sóc thú cưng. Hệ thống tập trung vào các đối tượng chính như khách hàng, thú cưng, bác sĩ thú y, nhân viên, tài khoản quản trị, lịch hẹn, hóa đơn, hồ sơ khám bệnh và phòng chăm sóc.

Trong quá trình thực hiện, đề tài áp dụng các kiến thức về phân tích yêu cầu, mô hình hóa dữ liệu, mô hình thực thể -- liên kết (ERD), ánh xạ ERD sang mô hình quan hệ và ngôn ngữ SQL. Cơ sở dữ liệu được xây dựng với các khóa chính, khóa ngoại và các ràng buộc toàn vẹn nhằm đảm bảo tính nhất quán của dữ liệu. Bên cạnh đó, dữ liệu mẫu và các truy vấn SQL được xây dựng để kiểm tra hoạt động của cơ sở dữ liệu.

Báo cáo trình bày toàn bộ quá trình từ phân tích nghiệp vụ, thiết kế ERD, ánh xạ sang mô hình quan hệ, cài đặt cơ sở dữ liệu, xây dựng dữ liệu mẫu, kiểm thử đến chuẩn hóa cơ sở dữ liệu. Qua đó, đề tài giúp vận dụng các kiến thức cơ sở dữ liệu vào một bài toán thực tế và tạo nền tảng cho việc mở rộng hệ thống trong tương lai.

Do kiến thức và thời gian thực hiện còn hạn chế, báo cáo không thể tránh khỏi những thiếu sót. Em  rất mong nhận được những ý kiến đóng góp và nhận xét để có thể tiếp tục hoàn thiện đề tài trong những phiên bản sau.

Em xin cảm ơn Thầy!
